\section{The finite metric evolution}\label{sec:finite-flow}

Let $\cP_r$ be the positive Hermitian determinant-one
matrices, with metric
$g_P(V,W)=\tr(P^{-1}VP^{-1}W)$. Its distance and geodesic are
\begin{equation}\label{eq:matrix-geometry}
 \begin{split}
 d(P,Q)^2&=\tr\!\left(\log(P^{-1/2}QP^{-1/2})^2\right),\\
 \Geo(P,Q;\theta)&=P^{1/2}
 \exp\!\left(\theta\log(P^{-1/2}QP^{-1/2})\right)P^{1/2}.
 \end{split}
\end{equation}
This is a complete simply connected symmetric space of
nonpositive curvature; see \cite[Sections~2.3--2.6]{bhatia}.
Unitary congruences act isometrically and fix $I$.
On $\cP_r^4$ use the product metric with weights $A_i$.
With physical weights \eqref{eq:weights}, define
\begin{equation}\label{eq:finite-potential}
 \Phi_{\tau,U}(\mathbf H)
 =\frac12\sum_{e=ij}w_e
          d(H_i,U_{ij}H_jU_{ij}^*)^2.
\end{equation}

\begin{proposition}\label{prop:finite-flow}
The negative gradient flow of \eqref{eq:finite-potential} is
\begin{equation}\label{eq:finite-flow}
 H_i^{-1}\dot H_i=\frac1{A_i}\sum_{e=ij}w_e B_{ij},
 \qquad B_{ij}=\log(H_i^{-1}U_{ij}H_jU_{ij}^*).
\end{equation}
For initial tuples in a fixed bounded ball, it exists for all
physical times $s\ge0$ and remains in a compact matrix range
independent of $\tau,U$. Distances between two such solutions
with the same parameters are nonincreasing. On that range the
vector field and every fixed-order tuple derivative are
$O(L_*^{-1})$, uniformly in all length ratios and $U$.
\end{proposition}
\begin{proof}
The gradient of $\frac12d(P,Q)^2$ with respect to $P$
is $-\Log_PQ=-P\log(P^{-1}Q)$. Summing the incident
terms and dividing by $A_i$ gives \eqref{eq:finite-flow}.
Each $B_{ij}$ is $H_i$-self-adjoint and trace free,
so the velocity is tangent to $\cP_r$. The reverse
identity is $B_{ji}=-U_{ij}^*B_{ij}U_{ij}$; no
commutation of the endpoint metrics is used.

Squared distance is jointly geodesically convex on
this Hadamard manifold. Hence $\Phi$ is convex and
vanishes at $\mathbf I=(I,I,I,I)$. If $D_A$ is the
weighted product distance, gradient-flow convexity gives
\begin{equation}\label{eq:finite-confinement}
 \frac d{ds}\frac12D_A(\mathbf H,\mathbf I)^2
 \le-\Phi(\mathbf H).
\end{equation}
Closed bounded subsets are compact, so the smooth
vector field continues for all time. Applying the
same convexity inequality at the two endpoints of
the joining geodesic proves nonexpansion between
two solutions. Finally \eqref{eq:finite-flow} is a
finite sum of smooth matrix functions with coefficients
$w_e\le2\pi/L_*$. All tuple derivatives are bounded
on the confined range, uniformly in the compact
unitary parameter space. In particular
$|\dot B_{ij}|\le C/L_*$.
\end{proof}

The zero set of $\Phi$ consists of compatible metric
tuples. On irreducible unitary representations this set
is $\{\mathbf I\}$ by Schur's lemma and determinant
one. On reducible representations it is generally
positive-dimensional.

\begin{lemma}\label{lem:finite-coercivity}
Put $S_i=\log H_i$ and $R_f=\sum_iA_i\|S_i\|^2$.
Then
\[
 \Phi(\mathbf H)\ge\tfrac12Q_{\tau,U}(S)
 \ge\tfrac12\nu_0R_f,\qquad
 R_f(s)\le e^{-\nu_0s}R_f(0).
\]
\end{lemma}
\begin{proof}
The exponential map at $I$ expands distances in this
nonpositively curved symmetric space. Since unitary
congruence commutes with the exponential,
$d(H_i,UH_jU^*)\ge\|S_i-US_jU^*\|$.
This proves the energy inequalities. Moreover
$R_f=D_A(\mathbf H,\mathbf I)^2$, so
\eqref{eq:finite-confinement} gives
$R_f'\le-2\Phi\le-\nu_0R_f$.
\end{proof}

The finite equation has the same linearization
$\dot X=-B_{\tau,U}X$ as the effective spectral
operator. Its nonlinear variables need not commute.
For Example~\ref{ex:noncommuting}, no fixed simultaneous
diagonalization of the cycle transports reduces the
equation to independent scalar vertex equations.
